\begin{statement}{C.2.UN.80}{Proved}{Every point-exceptional prime of an aligned parent is an actual boundary}
Let \(P\), \(\pi:P\to Z\), and \(q_Z:Z\to B\) be the
actual component and maps of C.2.UN.70. Write \(D_P\)
for the reduced union of its intersections with the
other central components. There is a nonzero constant
linear combination of the two restricted saturated
ambient translation fields which is the meromorphic
lift of a nowhere-zero holomorphic vertical field
\[
 W_Z\in H^0(Z,T_Z),\qquad dq_Z(W_Z)=0.
 \tag{C.EP.1}
\]
On every smooth elliptic fibre its complex-time period
lattice is the same fixed lattice. The field extends
without a zero or a pole across every smooth multiple
fibre of \(q_Z\).

Every prime exceptional for the specified point map
\(\pi\) is a polar divisor of its meromorphic lift
\(W_P\). Consequently
\[
                 \operatorname{Exc}(\pi)\subset D_P.
 \tag{C.EP.2}
\]
In particular, a last point-exceptional curve for
\(\pi\) is an actual boundary curve. This assertion
concerns the exceptional primes of this actual map:
it makes no assertion that every rational
\((-1)\)-curve on \(P\) is exceptional for \(\pi\).
The component and its ambient neighbourhood are not
changed.
\end{statement}

\begin{proof}
\par\medskip\noindent\textbf{The actual field space and its polar locus.}
Use the two-dimensional saturated field space of
C.2.UN.16 on the original multiplicity-one component
\(P\). In this full reduced situation no root cover
of \(P\) is needed. Denote a basis of its restrictions
by \(X,Y\). They are independent over \(\C\), commute,
and are meromorphic fields on \(P\).

Their actual ambient representatives are linear
combinations of holomorphic relative fields with
coefficients meromorphic in the disc coordinate.
They therefore have polar divisors only in the
central fibre. Saturation removes their pole along
the generic point of \(P\). On \(P\setminus D_P\)
the central fibre has just this one smooth branch,
so restriction gives
\[
       X,Y\ \hbox{holomorphic on }P\setminus D_P .
 \tag{C.EP.3}
\]
This is the exact restriction comparison used in
(C.MR.8) and C.2.UN.70. It applies to every constant
linear combination. All boundary and exceptional
curves of the specified map \(\pi\) are vertical
for \(q=q_Z\pi\).

Remove finitely many values from \(B\) so that
\(\pi\) is an isomorphism over their complement
\(B^\circ\), all fibres are smooth, and the fields
and their projected coefficients are regular.
On a compact elliptic fibre the latter coefficients
are holomorphic functions and hence constants.
Thus \(X,Y\) project to meromorphic vector fields
on \(B\), exactly as in the alignment proof.

\par\medskip\noindent\textbf{Rank one: the period identity gives a constant combination.}
Suppose the meromorphic rank of \(\langle X,Y\rangle\)
is one. Alignment identifies its connected elliptic
first integral with the actual map \(q\). Thus
\[
                       Y=(q^*H)X
 \tag{C.EP.4}
\]
for a nonconstant meromorphic function \(H\) on
\(B\). The actual return-germ calculation
(C.RG.4)--(C.RG.12), performed on the original
ambient representatives, gives time periods
\(\lambda_1,\lambda_2\) of \(X\) and constants
\(d_i,k_i\in\C\) such that
\[
            \lambda_i=d_i-k_iH,\qquad
            \lambda_2/\lambda_1=\tau\in\mathbb H .
 \tag{C.EP.5}
\]
The formula continues on \(B^\circ\) with these
same constants. This uses the actual global ratio
\(H\), the continued integral cycles, and the
identity theorem; it does not extend \(q\) to a
threefold neighbourhood.

Take
\[
                       W=d_1X-k_1Y=\lambda_1X.
 \tag{C.EP.6}
\]
A nonzero integral cycle on an elliptic curve has
nonzero time period, so \(W\) is nonzero on every
fibre under consideration. Multiplying a field by
\(\lambda_1\) divides its time periods by
\(\lambda_1\); the two periods of \(W\) are exactly
\(1,\tau\). Thus (C.EP.6) is a constant linear
combination with a fixed nonzero period lattice.
In an integration coordinate on each fibre it
is the fixed invariant translation field.

\par\medskip\noindent\textbf{Rank two: commuting transverse transport preserves the periods.}
Suppose instead that the meromorphic rank is two.
The two projected fields on the curve \(B\)
commute. Two nonzero commuting meromorphic fields
on a connected curve have constant ratio:
writing them locally as \(a(z)\partial_z\) and
\(b(z)\partial_z\), the bracket equation is
\(ab'-ba'=0\). The ratio therefore continues as
a constant. They cannot both project to zero,
since vertical fields on an elliptic fibre have
meromorphic rank at most one. After interchanging
the basis if necessary there is consequently a
constant \(c\) for which
\[
             W=Y-cX\ne0,\qquad q_*W=0,\qquad
                         q_*X\ne0 .
 \tag{C.EP.7}
\]
On the compact fibres over \(B^\circ\), \(W\)
is a nonzero translation field. The local flow
of the transverse field \(X\), constructed on
an entire compact fibre for sufficiently small
time, identifies neighbouring elliptic fibres.
Since \([X,W]=0\), this identification preserves
the time form and its two integral periods.
Thus their lattice is locally constant and
hence constant on connected \(B^\circ\).
This is the actual compact-fibre flow comparison
of (C.RF.1)--(C.RF.2). A constant rescaling of
\(W\) makes one nonzero period equal to \(1\).
Once again \(W\) is a constant combination of
the original two restricted fields.

\par\medskip\noindent\textbf{Extension on the specified minimal elliptic model.}
On \(q_Z^{-1}(B^\circ)\), push \(W\) through
the open isomorphism of \(\pi\). Integration of
its fixed lattice identifies each small smooth
fibre tube with a torsor under one fixed elliptic
curve \(F\). Every overlap has translation
linear part, so its derivative preserves the
same invariant field of \(F\).

At a multiple value, use the actual normalized
root chart of C.2.UN.70 and (C.RG.13):
\[
 (F\times\Delta_z)/\langle
       (u,z)\longmapsto(u+\eta,\zeta_mz)\rangle,
       \qquad t=z^m .
 \tag{C.EP.8}
\]
The translation \(\eta\) gives the free cyclic
action on the total space. The invariant
vertical field on \(F\times\Delta_z\) is fixed
by this action and is nowhere zero. On the
punctured root disc our field agrees with it:
its coefficient relative to the invariant
frame has the fixed nonzero time periods and
therefore is the same constant. It follows
that the agreement extends across \(z=0\)
and descends through the actual etale quotient.
These extensions agree on overlaps with the
previously fixed field. This proves (C.EP.1),
including every zero or pole that might have
been hidden at an omitted base value.

The meromorphic lift of \(W_Z\) through \(\pi\)
agrees with the constant combination (C.EP.6)
or (C.EP.7) on their common open isomorphism,
and therefore agrees everywhere as a
meromorphic vector field. In particular its
polar divisors are contained in \(D_P\), by
(C.EP.3).

\par\medskip\noindent\textbf{Every exceptional valuation is polar.}
Let \(E\) be any prime exceptional for \(\pi\),
with image \(p\in Z\). The nowhere-zero field
\(W_Z\) has a holomorphic flow box \((z,w)\)
centred at \(p\) in which
\[
                   W_Z(z)=1,\qquad W_Z(w)=0.
 \tag{C.EP.9}
\]
Both pullbacks vanish to positive order along
\(E\). Work at a generic smooth point of \(E\)
and let \(s\) be its local equation. Write
\[
            \pi^*w=s^b u,\qquad
            \pi^*z=s^a v,\qquad a,b\ge1,
 \tag{C.EP.10}
\]
where \(u,v\) are holomorphic units after
shrinking this generic neighbourhood.

Suppose the lifted derivation \(D=W_P\) were
holomorphic there. The identity
\(D(\pi^*w)=0\), divided by \(s^{b-1}\) and
reduced modulo \(s\), gives
\[
                    b\,u\,D(s)=0\pmod{s}.
 \tag{C.EP.11}
\]
Characteristic zero and the fact that \(u\)
is a unit imply \(D(s)\in(s)\). Thus \(D\)
preserves the ideal of \(E\), and applying it
to the second expression in (C.EP.10) gives
\(D(\pi^*z)\in(s)\). This contradicts the
exact pulled-back identity \(D(\pi^*z)=1\).
Hence \(E\) is a polar divisor. The argument
applies directly to every exceptional prime,
including all infinitely-near point divisors;
it does not assume that a first pole survives
an unexamined modification.

Combining this with (C.EP.3) proves (C.EP.2).
The exact period rescalings, the root-chart
invariance, and finite toric valuation
certificates for (C.EP.9)--(C.EP.11) are
executed in \path{code/A_translation_exceptions.py}.
The general exceptional-prime assertion is
the local derivation proof above, without a
toric hypothesis on the actual point map.
\end{proof}


% C.2.UN.65 from part_C2_trivial_canonical.tex

